
La structure Hilbertienne permet de définir la projection sur un convexe fermé, et d'identifier l'espace avec son dual via le théorème de Riesz, voir les \cref{th:projConvexe,th:reprRiesz}. Ces résultats puissants trivialisent pour l'essentiel les questions de réflexivité et de séparation de convexes étudiées dans le cas des espaces de Banach, voir les \cref{subsec:reflexivite,subsec:HahnBanachGeom}.

Dans cette section, $\bK$ désigne le corps $\bR$ ou $\bC$.

\begin{retenir}
	L'inégalité de Cauchy-Schwarz, la définition et la caractérisation de la projection sur un convexe fermé, le théorème de Représentation de Riesz, de la base Hilbertienne.
	(\Cref{prop:CauchySchwarz,th:projConvexe,th:reprRiesz,th:baseHilbert}.)
\end{retenir}

\inventaire{
\Nref{def:prehilbert} \\
\Nref{prop:CauchySchwarz} \\
\Nref{def:Hilbert} \\
\Nref{th:projConvexe} \\
\Nref{th:reprRiesz} \\
\Nref{ex:l2Hilbert} (Hilbert $L^2(X,\mu)$)\\
\Nref{prop:decompOrth} \\
\Nref{prop:decompCone} \\
\Cref{ex:coneSdp}(Cône dual de $S_d^+$) \\
\Nref{th:baseHilbert}
}

\begin{definition}[Espace préhilbertien]
\label{def:prehilbert}
	Une forme sesquilinéaire sur un $\bK$-espace vectoriel $E$ est une application $\vp : E \times E \to \bK$ qui satisfait pour tous $x,x',y,y' \in E$, $\lambda \in \bK$ : 
	\begin{itemize}
		\item (Linéarité à droite) $\vp(x,\lambda y) = \lambda \vp(x,y)$, $\vp(x,y+y')=\vp(x,y)+\vp(x,y')$. %, $\forall x,y,y' \in E$, $\lambda \in \bK$
		\item (Antilinéarité à gauche) $\vp(\lambda x,y) = \overline \lambda \vp(x,y)$, $\vp(x+x',y)=\vp(x,y)+\vp(x',y)$.  
	\end{itemize}
	Elle est dite 
	\begin{itemize}
		\item \emph{symétrique} si $\vp(x,y) = \overline {\vp(y,x)}$, $\forall x,y \in E$. 
		\item \emph{positive} si $\vp(x,x) \in \bR_+$, $\forall x \in E$.
		\item \emph{définie} si $\vp(x,x) \neq 0$, $\forall x \in E\sm \{0\}$.
	\end{itemize}
	Un \emph{produit scalaire} $\<\cdot,\cdot\>$ est une forme sesquilinéaire symétrique définie positive. Dans ce cas, l'espace est dit \emph{préhilbertien}, et on note $\|x\| := \sqrt{\<x,x\>}$ pour tout $x \in E$. 
\end{definition}

Tous les éléments $x,y$ d'un espace préhilbertien, satisfont les égalités suivantes : 
\begin{align*}	
	\|x+y\|^2 &= \|x\|^2 + \|y\|^2 + 2 \Re(\<x,y\>), &
	\|x+y\|^2 + \|x-y\|^2 &= 2 (\|x\|^2 + \|y\|^2),
\end{align*}
la seconde étant appelée identité du parallélogramme. 

\begin{proposition}[Inégalité de Cauchy-Schwarz]
\label{prop:CauchySchwarz}
Soit $(E,\<\cdot ,\cdot\>)$ un espace préhilbertien. Alors pour tous $x,y \in E$, 
\begin{equation}
\label{eq:CauchySchwarz}
	|\<x,y\>| \leq \|x\| \|y\|, 
\end{equation}
et $\|\cdot\|$ définit une norme sur $E$. L'égalité a lieu ssi $x$ et $y$ sont colinéaires. 
\end{proposition}

\begin{proof}
	L'égalité a clairement lieu dans \eqref{eq:CauchySchwarz} si $x$ et $y$ sont colinéaires. On  suppose désormais qu'ils ne le sont pas. 
	Soit $\alpha\in \bK$, $|\alpha|=1$, tel que $\Re(\alpha\<x,y\>) = |\<x,y\>|$.
	Considérons le polynôme %du second degré 
	\begin{equation*}
		P(t) := \|x+\alpha t y\|^2 
		= \<x+\alpha t y, x+\alpha t y\>
		= \<x,x\> + 2 t \Re(\alpha\<x,y\>)+ t^2\<y,y\> 
		= \|x\|^2 + 2t|\<x,y\>| + t^2\|y\|^2,
	\end{equation*}
	d'indéterminée $t \in \bR$.
	Comme $P$ est du second degré et ne s'annule pas, son discriminant est strictement négatif, d'où \eqref{eq:CauchySchwarz}. 
	
	La fonction $\|\cdot\|$ est clairement définie, positive et homogène, conformément aux axiomes de la \nref{def:norme}. Il reste donc à vérifier l'inégalité triangulaire : $\|x+y\|^2 = \|x\|^2 + 2\Re(\<x,y\>) + \|y\|^2 \leq \|x\|^2 + 2 \|x\| \|y\| + \|y\|^2 = (\|x\|+\|y\|)^2$, $\forall x,y \in E$, ce qui conclut. 
\end{proof}

\begin{definition}[Hilbert]
\label{def:Hilbert}
%Un espace vectoriel normé dont la norme découle d'un produit scalaire préhilbertien. 
Un espace préhilbertien complet est un Hilbert. 
\end{definition}

\begin{theorem}[Projection sur un convexe fermé]
\label{th:projConvexe}
Soit $H$ Hilbert et $K\subset H$ un convexe fermé non-vide. 
Pour tout $x \in H$, le problème suivant admet une unique solution
\begin{equation}
\label{eqdef:proj_cvx_ferme}
	P_K(x) := \argmin_{y \in K} \|x-y\|,
\end{equation}
définissant la projection $P_K : H \to K$. Elle est caractérisée par la condition d'angle obtus : pour tout $y \in K$
\begin{equation}
\label{eq:proj_obtus}
	\Re\<x-P_K(x), y-P_K(x)\> \leq 0.
\end{equation}
Enfin, $P_K$ est $1$-Lipschitz, et plus précisément pour tous $x,y \in H$
\begin{equation}
\label{eq:proj_moyennee}
	\|P_K(x)-P_K(y)\|^2 \leq \Re\<x-y,\ P_K(x)-P_K(y)\>.
\end{equation} 
\end{theorem}

\begin{proof}
Notons $\delta:=d_K(x)$ l'infimum de \eqref{eqdef:proj_cvx_ferme}, et 
soit $(x_n)_{n \in \bN}$ une suite minimisante, i.e.\ $x_n\in K$ et $\ve_n := \|x_n - x\|-d_K(x) \to 0$ lorsque $n \to \infty$.
Etant donnés $p,q \geq N$ on a 
\begin{equation*}
	\delta^2 \leq \big\| x - \frac{x_p+x_q} 2\big\|^2 
	= \frac 1 2 \big(\|x-x_p\|^2 + \|x-x_q\|^2\big) - \frac 1 4\|x_p-x_q\|^2
	\leq \frac 1 2 \big((\delta+\ve_p)^2 + (\delta+\ve_q)^2\big)	- \frac 1 4\|x_p-x_q\|^2,
\end{equation*}
par l'identité du parallélogramme, d'où $\|x_p-x_q\| \leq 2\sqrt{\delta(\ve_p+\ve_q)+ \frac 1 2 (\ve_p^2+\ve_q^2)}$. Ainsi $(x_n)_{n \in \bN}$ est de Cauchy, donc convergente par complétude de $H$, et par continuité de la norme et fermeture de $K$ sa limite est un minimiseur de \eqref{eqdef:proj_cvx_ferme}.
On déduit aussi de cette preuve l'unicité du minimiseur de \eqref{eqdef:proj_cvx_ferme}.

Considérons $x_* := P_K(x)$, et $y\in K$ arbitraire. On a $(1-t) x_*+ ty \in K$ pour tout $t \in [0,1]$, par convexité de $K$, puis
\begin{equation}
\label{eq:caract_proj_proof}
	\|x_*-x\|^2 \leq \|(1-t) x_* + t y - x\|^2 = \|x-x_*\|^2 - 2 t \Re(\<x-x_*,y-x_*\>) + t^2 \|y-x_*\|^2
\end{equation}
par minimalité de $x_*$ dans \eqref{eqdef:proj_cvx_ferme}. On en déduit $\Re(\<x-x_*,y-x_*\>)\leq 0$ par un développement limité en $t=0^+$, comme annoncé. 
Réciproquement, considérons $x_* \in K$ tel que $\Re(\<x-x_*,y-x_*\>)\leq 0$ pour tout $y \in K$.
Alors $\|y-x\|^2 = \|x-x_*\|^2 - 2 \Re(\<x-x_*,y-x_*\>) + \|y-x_*\|^2 \geq \|x-x_*\|^2$ pour tout $y \in K$, donc $x_*$ minimise \eqref{eqdef:proj_cvx_ferme}, comme annoncé.

Finalement \eqref{eq:proj_moyennee} s'obtient en sommant les inégalités $\Re\<x-P_K(x), P_K(y)-P_K(x)\> \leq 0$ et $\Re\<y-P_K(y), P_K(x)-P_K(y)\> \leq 0$, issues de \eqref{eq:proj_obtus}. La Lipschitzianité de $P_K$ en découle par la \nref{prop:CauchySchwarz}.
\end{proof}

La propriété \eqref{eq:proj_moyennee} est équivalente à la $1$-Lipschitzianité de $Q := 2P_K-\Id$. On dit que $P_K := (\Id+Q)/2$ est un \emph{opérateur moyenné}: c'est la moyenne de l'identité et d'un opérateur $1$-Lipschitz.


Le résultat suivant montre que l'on peut identifier un Hilbert à son dual. Les normes se correspondent isométri\-quement, ainsi que les topologies faible et *faible. 

\begin{theorem}[Représentation de Riesz]
\label{th:reprRiesz}
Soit $H$ un espace de Hilbert. Alors $H \to H^* : x \mapsto \vp_x := \<x,\cdot\>$, est une bijection isométrique anti-linéaire.  
\end{theorem}

\begin{proof}
	\emph{Preuve de l'isométrie}. On a $\|\vp_x\|_{H^*} \geq \vp_x( x/\|x\|) = \|x\|$ pour tout $x \neq 0$, et par ailleurs $|\vp_x(y)| = |\<x,y\>| \leq \|x\| \|y\|$ par Cauchy-Schwarz, d'où $\|\vp_x\|_{H^*} = \|x\|$, comme annoncé.
	
	\emph{Preuve de la surjectivité}. Soit $\vp \in H^*\sm \{0\}$, et soit $K := \vp^{-1}\{0\}$ son noyau, qui est convexe en tant qu'hyperplan, et fermé car $\vp$ est continue. 
	Soit $x_0\in H$ tel que $\vp(x_0)=1$, qui existe par homogénéité, et soit $x_* := x_0 - P_K(x_0)$, de sorte que $\vp(x_*) = \vp(x_0) = 1$ et $\<x_*,y\> = 0$ pour tout $y \in K$. Alors pour tout $x\in H$ on a $\vp(x-\vp(x) x_*) = 0$ par linéarité, donc $x-\vp(x) x_* \in K$, donc $\<x_*,x-\vp(x) x_*\> = 0$, donc $\vp(x) = \<x_*,x\>/\|x_*\|^2$ ce qui conclut. 
\end{proof}

\begin{exemple}
\label{ex:l2Hilbert}
	Soit $(X,\mu)$ un espace mesuré. Alors $L^2(X,\mu)$ est un Hilbert, muni du produit scalaire $\<f,g\> := \int_X \overline f(x) g(x) \diff\mu(x)$, voir l'\cref{ex:LpBanach}. En particulier, toute forme linéaire continue s'y écrit $g \mapsto \int_X \overline f(x) g(x) \diff\mu(x)$, pour un certain $f \in L^2(X, \mu)$.
\end{exemple}

%Les résultats suivants illustrent les décompositions d'élements d'un Hilbert obtenues grâce 
La projection orthogonale permet de décomposer les éléments d'un Hilbert.

\begin{proposition}[Décomposition orthogonale]
\label{prop:decompOrth}
	Soit $H$ Hilbert. 
	La projection sur un sous-espace vectoriel fermé $F \subset H$ est caractérisée par la condition $\<x-P_F(x),y\> = 0$, $\forall y \in F$.
	De plus $F \oplus F^\perp = H$, $P_F+P_{F^\perp} = \Id$, et $P_F$ et $P_{F^\perp}$ sont linéaires continues, où $F^\perp := \{y \in H \mid \forall x \in F,\, \<x,y\>=0\}$ désigne le sous-espace orthogonal.
\end{proposition}

\begin{proof}
	Soit $x \in H$ et $p := P_F(x) \in F$. D'après \eqref{eq:proj_obtus} on a $\Re\<x-p,y-p\> \leq 0$ pour tout $y \in F$. Donc $\Re\<x-p,y\> = 0$ pour tout $y \in F$, par stabilité additive de $F$, et donc $\<x-p,y\> = 0$, par stabilité de $F$ par multiplication par un scalaire, comme annoncé. La linéarité de cette caractérisation de $P_F(x)$ implique la linéarité du projecteur $P_F$, qui est également continu par le \cref{th:projConvexe}, comme annoncé.
	
	Notons $z := x-p$, et observons que $z \in F^\perp$ par la discussion précédente. Par ailleurs $x-z = p \in F \subset (F^\perp)^\perp$, donc par la caractérisation \eqref{eq:proj_obtus} on a $z = P_{F^\perp} x$. Ainsi $x = P_F x + P_{F^\perp} x$, comme annoncé, et ainsi $F + F^\perp = H$. Remarquons finalement que $F \cap F^\perp = \{0\}$ (si $x \in F \cap F^\perp$, alors $\|x\|^2 = \<x,x\> = 0$), donc $F \oplus F^\perp = H$ comme annoncé, ce qui conclut.
\end{proof}

Une variante du résultat classique précédent traite des cônes convexes fermés. Rappelons qu'un cône convexe est une partie $K$ d'un espace $\bR$-vectoriel telle que $\lambda x+\mu y \in K$ pour tous $x,y \in K$ et tous $\lambda,\mu \geq 0$. 

\begin{proposition}[Décomposition orthogonale associée à un cône \facultatif]
\label{prop:decompCone}
	Soit $H$ un $\bR$-Hilbert, soit $K \subset H$ un cône convexe fermé, et soit $K^\circ := \{y\in H\mid \forall x \in K,\ \<x,y\> \leq 0\}$ le cône polaire. %; noter que $K \cap K^\circ = \{0\}$.
	On a la caractérisation suivante des projections sur $K$ et $K^\circ$ : %sont caractérisées par %$P_K+P_K^\perp = \Id$ et $\<P_K,P_{K^\circ}\> = 0$, 
	\begin{equation}
	\label{eq:proj_cone}
		(y = P_K(x),\ z=P_{K^\circ} (x)) \ \Leftrightarrow \ (y\in K,\, z \in K^\circ,\, \<y,z\>=0 \text{ et } x = y + z). 
	\end{equation}
\end{proposition}

\begin{proof}
	Supposons (\ref{eq:proj_cone}, droite). Alors pour tout $y' \in K$ on a $\<x-y,y'-y\> = \<z,y'-y\> = \<z,y'\> \leq 0$, donc $y = P_K(x)$ par la caractérisation \eqref{eq:proj_obtus}. De même $z = P_{K^\circ} (x)$, comme annoncé.
	
	Réciproquement, considérons $y:=P_K(x)$ et posons $z := x-y$. On a $\<x-y,y'-y\> =\<z,y'-y\> \leq 0$, $\forall y' \in K$, par \eqref{eq:proj_obtus}. Par la propriété de cône on obtient $t \<z,y'\> \leq \<z,y\>$ pour tout $t \geq 0$, donc $\<z,y'\> \leq 0$, et donc $z \in K^\circ$. 
	A l'inverse, en choisissant $y'=0$ on obtient $\<z,0-y\>\leq 0$, donc $\<z,y\>=0$. Finalement $\<x-z,z'-z\> = \<y,z'-z\> = \<y,z'\> \leq 0$ pour tout $z' \in K^\circ$, donc $z = P_{K^\circ}(x)$, ce qui conclut.  
\end{proof}

\begin{exemple}[Cône dual de $\bR_+^d$ et $\cS^+_d$ \facultatif]
\label{ex:coneSdp}
	On note parfois $K^* := - K^\circ  = \{y \in H\mid \forall x \in K, \<x,y\> \geq 0\}$ le cône dual, d'un cône $K$ donné. 
	Un cône est dit auto-dual si $K = K^*$. Par exemple $\bR_+^d \subset \bR^d$ et $\cS_d^+ \subset \cS_d$ sont des cônes auto-duaux. La décomposition de la \cref{prop:decompCone} associée au cône $\bR_+^d$ correspond aux parties positive et négative de $x=(x_i)_{1 \leq i \leq d} \in \bR^d$~: 
	\begin{align*}
		x &= x^+ - x^-, &\text{où } x^+_i &:= \max\{0,x_i\},\ x^-_i := \max\{0,-x_i\}
	\end{align*}
	De même, la décomposition de la \cref{prop:decompCone} associée au cône $\cS_+^d$ correspond aux parties positive et négative du spectre : pour $D \in \cS_d$, diagonalisée sous la forme $U^\top \Diag(\lambda_1,\cdots,\lambda_d) U$ avec $U$ une matrice orthogonale, on a 
	\begin{equation*}
		D  
		= U^\top \Diag(\max \{\lambda_1,0\},\cdots,\max \{\lambda_d,0\}) U - U^\top \Diag(\max \{-\lambda_1,0\},\cdots,\max \{-\lambda_d,0\}) U.
	\end{equation*}
\end{exemple}

%En optimisation, le cône normal est le dual du cône intérieur, en un point $x$ d'un convexe $K$.
%\begin{align*}
%	\cC_K^\mathrm{int}(x) &= \{v \in H \mid \exists t >0, x+t v\in K\},&
%	\cC_K^\mathrm{norm}(x) &= \{v \in H \mid \forall y \in K, \<v,y-x\> \leq 0\}.
%\end{align*}
%Noter que le cône normal est fermé, mais pas forcément le cône intérieur.
%Dans un Banach, des notions analogues sont discutées après \nref{def:convexePolaire}.

%\todoi{
%Si nécessaire, évoquer le fait que les Hilbert séparables de dimension infinie sont isomorphes à $L^2(\bN,\bK)$.
%}

La structure Hilbertienne est si forte que $\ell^2 := L^2(\bN,\bK)$ est, à isométrie linéaire près, son seul représentant intéressant, comme le montre le théorème suivant qui conclut cette introduction. 

\begin{theorem}[Base Hilbertienne]
\label{th:baseHilbert}
Soit $H$ un espace de Hilbert. Sont équivalents : 
\begin{enumerate}[(i)]
	\item $H$ est séparable et de dimension infinie.
	\item $H$ admet une \emph{base hilbertienne}, c'est-à-dire une famille $(e_n)_{n\in \bN}$ de vecteurs unitaires deux à deux orthogonaux, telle que $\Vect\{e_n \mid n \in \bN\}$ est dense dans $H$.
	\item Il existe une bijection linéaire et isométrique entre $H$ et $\ell^2 := L^2(\bN,\bK)$.
\end{enumerate}
En supposant (ii), la bijection de (iii) est donnée par $x \in H \mapsto (\<e_n,x\>)_{n \in \bN} \in \ell^2$, et sa réciproque par $(x_n)_{n \in \bN} \mapsto \sum_{n \in \bN} x_n e_n$.
\end{theorem}

\begin{proof}
	\emph{Preuve de (i) $\Rightarrow$ (ii).} Soit $(x_n)_{n \in \bN}$ une suite dense dans $H$, et soit $F_N := \Vect\{x_n \mid n < N\}$ pour tout $N \geq 0$, avec la convention $F_0:=\{0\}$.
	Comme $F_n \subset F_{n+1}$ pour tout $n \geq 0$, comme l'espace vectoriel $\cup_{n \in \bN} F_n$ est dense dans $H$ par hypothèse, et comme $H$ est de dimension infinie, on a $\dim(F_n) \to \infty$ lorsque $n \to \infty$; de plus, $\dim(F_{n+1}) \leq \dim(F_n)+1$ pour tout $n \in \bN$ par construction. 
	Ainsi, il existe $\vp : \bN \to \bN$ extractrice telle que $\dim(F_{\vp(n)}) = n$ pour tout $n \in \bN$; quitte à extraire, on peut supposer $\dim(F_n) = n$. 
	On vérifie alors facilement que $e_n := (x_n - P_{F_n}(x_n)) / \|x_n - P_{F_n}(x_n)\|$, pour tout $n \in \bN$, définit une base Hilbertienne.
	
	\emph{Preuve de (ii) $\Rightarrow$ (iii).}
	La projection orthogonale de $x \in H$ sur $F_N$ est donnée par $P_{F_N}(x) = \sum_{n< N} \<e_n,x\>e_n$, pour tout $N \geq 0$, d'après le critère de la \cref{prop:decompOrth}, et donc $\|P_{F_N}(x)\|^2 = \sum_{n<N} |\<e_n,x\>|^2$. 
	Ainsi $\|x\|^2 - \sum_{n<N} |\<e_n,x\>|^2 = \|x-P_{F_N}(x)\|^2 \to 0$ lorsque $N \to \infty$, ce qui montre la bonne définition et l'isométrie.
	Réciproquement, si $(x_n)_{n\in \bN} \in \ell^2$, alors les sommes partielles $s_N := \sum_{n<N} x_n e_n$ satisfont le critère de Cauchy, puisque $\|s_n-s_m\|^2 = \sum_{n\leq k <m} |x_k|^2$, donc convergent vers un élément $x\in H$ par complétude, et celui-ci satisfait $\<e_n,x\> = \lim_{N \to \infty} \<e_n,s_N\> = x_n$ pour tout $n \in \bN$ par continuité du produit scalaire.
	
	\emph{Preuve de (iii) $\Rightarrow$ (i).} Les éléments de $\ell^2$ à support fini et coefficients dans $\bQ$ forment une famille dénombrable et dense, d'où la séparabilité. De plus $\ell^2$ contient la famille libre infinie  $(\delta^N)_{N\geq 0}$, sa base hilbertienne canonique, définie par $\delta^N_n = 1$ si $n=N$ et $\delta^N_n = 0$ sinon.
\end{proof}

